\section{Introducción}

\textbf{Planteamiento del problema.} \\
El reto en esta primera actividad es entender cómo pasar operaciones matemáticas básicas a un lenguaje de programación real como es C. Específicamente, necesitamos ver cómo declarar variables de tipo entero para guardar valores fijos y aplicar los operadores correctos para sumar, multiplicar y obtener el residuo de una división[cite: 3].

\vspace{0.3cm}

\textbf{Motivación.} \\
Hacer este ejercicio nos sirve para aterrizar las bases de la programación estructurada. Dominar la sintaxis inicial —como el uso de las librerías, la función principal y el operador de módulo— es lo que nos permitirá construir programas más complejos después[cite: 4]. Es, básicamente, asegurar que los cimientos de nuestra lógica de programación estén bien puestos.

\vspace{0.3cm}

\textbf{Objetivos.} \\
Lo que buscamos con este programa es familiarizarnos con la estructura básica de un archivo en C y verificar que nuestro entorno de trabajo funcione bien. El objetivo principal es practicar la declaración de variables y el uso de operadores aritméticos básicos para procesar datos de forma organizada[cite: 5].